\chapter{Réseaux de neurones pour les données tabulaires bruitées}\label{ch:ml:neural-networks-for-noisy-tabular-data}

Deux chercheurs entraînent le même petit réseau, sur les mêmes données, avec le même code, et rapportent des coefficients
d'information de 0.057 et 0.065 sur les mêmes dix années de test. Ils discutent de l'utilité d'une troisième couche cachée~; en moyenne, elle ajoute
0.002. La seule différence entre leurs deux exécutions est la graine aléatoire, et elle vaut quatre fois l'architecture
sur laquelle ils se disputent. Sur les données de marché, un \omterm{def:ml:neural-networks-for-noisy-tabular-data:mlp}{réseau de neurones} est une variable aléatoire avant d'être un modèle~: ses poids initiaux,
ses masques de \omterm{def:ml:neural-networks-for-noisy-tabular-data:dropout}{dropout} et l'ordre de ses \omterm{def:ml:neural-networks-for-noisy-tabular-data:backprop}{mini-lots} font tous bouger le résultat, et le signal est trop faible pour les noyer.
Ce chapitre construit le \omterm{def:ml:neural-networks-for-noisy-tabular-data:mlp}{perceptron multicouche}, l'entraîne comme l'exigent les données tabulaires bruitées, et mesure ses graines comme une
source de variance.

\section{Le perceptron multicouche et la rétropropagation}

\begin{definition}[Réseau de neurones, perceptron multicouche, fonction d'activation]\label{def:ml:neural-networks-for-noisy-tabular-data:mlp}
Un \emph{réseau de neurones}\index{réseau de neurones} (neural network) est une fonction construite par composition d'applications linéaires paramétrées avec
des fonctions non linéaires fixées. Un \emph{perceptron multicouche}\index{perceptron multicouche} (multilayer perceptron, MLP) est le cas entièrement connecté~:
$z_0 = x$, $z_l = \mathrm{act}(A_lz_{l-1} + c_l)$ pour $l = 1,\dots,L$, et $f_\theta(x) = A_{L+1}z_L + c_{L+1}$, avec
paramètres $\theta = (A_l, c_l)_l$. La \emph{fonction d'activation}\index{fonction d'activation} (activation function) $\mathrm{act}$ s'applique
coordonnée par coordonnée~; la plus courante est $\mathrm{relu}(u) = \max(u, 0)$.
\end{definition}

\begin{omfigure}
\begin{tikzpicture}[x=2.4cm,y=0.55cm,font=\footnotesize,
  neuron/.style={circle,draw=omDef,fill=omDef!12,minimum size=9pt,inner sep=0pt}]
\foreach \i in {1,...,6} { \node[neuron] (a\i) at (0,\i-3.5) {}; }
\foreach \i in {1,...,5} { \node[neuron] (b\i) at (1,\i-3) {}; }
\foreach \i in {1,...,3} { \node[neuron] (c\i) at (2,\i-2) {}; }
\node[neuron,draw=omThm,fill=omThm!15] (d) at (3,0) {};
\foreach \i in {1,...,6} { \foreach \j in {1,...,5} { \draw[gray!45] (a\i) -- (b\j); } }
\foreach \i in {1,...,5} { \foreach \j in {1,...,3} { \draw[gray!45] (b\i) -- (c\j); } }
\foreach \i in {1,...,3} { \draw[gray!45] (c\i) -- (d); }
\node[align=center] at (0,-4.6) {40 caractéristiques\\classées};
\node[align=center] at (1,-4.6) {32 unités\\$\mathrm{relu}(A_1x + c_1)$};
\node[align=center] at (2,-4.6) {16 unités\\$\mathrm{relu}(A_2z_1 + c_2)$};
\node[align=center] at (3,-4.6) {prévision\\$A_3z_2 + c_3$};
\end{tikzpicture}

\omcaption{Le petit réseau du chapitre (dessiné avec moins d'unités qu'il n'en a)~: 40 entrées, couches cachées de 32 et 16
unités, une sortie, environ 1\,850 paramètres. Le réseau plus grand ajoute une première couche de 64 unités.}\label{fig:ml:nets:mlp}
\end{omfigure}

\begin{definition}[Rétropropagation, mini-lot, époque]\label{def:ml:neural-networks-for-noisy-tabular-data:backprop}
La \emph{rétropropagation}\index{rétropropagation} (backpropagation) calcule le gradient de la perte par rapport à chaque paramètre d'un
réseau par différentiation algorithmique en mode inverse (livre~4, chapitre~28)~: une passe avant mémorise les valeurs des
couches, puis une passe arrière propage $\partial\mathcal L/\partial z_l$ de la sortie vers l'entrée. Un \emph{mini-lot}\index{mini-lot} (mini-batch)
est le sous-ensemble aléatoire de lignes d'entraînement sur lequel on calcule un pas de gradient~; une \emph{époque}\index{epoque@époque} (epoch) est un passage
complet sur l'\omterm{def:ml:why-financial-machine-learning-is-different:sets}{ensemble d'entraînement}, découpé en mini-lots.
\end{definition}

Le coût du gradient est un petit multiple du coût de la passe avant, quel que soit le nombre de paramètres~: c'est le
résultat du livre~4 sur le mode inverse, et c'est pourquoi les réseaux peuvent compter des millions de paramètres. Les réseaux du chapitre en ont entre deux et
cinq mille et s'entraînent en une seconde environ sur un seul fil d'exécution CPU.

\section{Optimiseurs, normalisation et régularisation}

\begin{definition}[Adam, décroissance des poids]\label{def:ml:neural-networks-for-noisy-tabular-data:adam}
\emph{Adam}\index{Adam} est une descente de gradient stochastique (livre~4, chapitre~24) avec un pas propre à chaque paramètre~: avec le gradient
$g_k$ à l'étape $k$, il conserve des moyennes $m_k = \beta_1m_{k-1} + (1-\beta_1)g_k$ et $v_k = \beta_2v_{k-1} + (1-\beta_2)g_k^2$
et déplace $\theta$ de $-\eta\,\hat m_k/(\sqrt{\hat v_k} + \epsilon)$, où les chapeaux corrigent le biais des moyennes vers zéro
(Kingma et Ba). La \emph{décroissance des poids}\index{décroissance des poids} (weight decay) réduit chaque poids de $\eta\lambda\theta$ à chaque pas, en dehors
du pas de gradient (AdamW, Loshchilov et Hutter)~; pour une descente de gradient ordinaire, c'est une pénalité ridge.
\end{definition}

\begin{definition}[Dropout, normalisation par lot et par couche]\label{def:ml:neural-networks-for-noisy-tabular-data:dropout}
Le \emph{dropout}\index{dropout} met à zéro chaque unité cachée avec probabilité $d$ pendant l'entraînement, en multipliant les autres par
$1/(1 - d)$, et utilise le réseau complet à la prédiction. La \emph{normalisation par lot}\index{normalisation par lot} (batch normalisation) standardise
chaque unité cachée sur les lignes du \omterm{def:ml:neural-networks-for-noisy-tabular-data:backprop}{mini-lot} (moyennes glissantes à la prédiction)~; la \emph{normalisation par
couche}\index{normalisation par couche} (layer normalisation) standardise les unités de chaque ligne, indépendamment du lot.
\end{definition}

Les données du chapitre sont le panel du chapitre~4 (500 actions, quarante caractéristiques classées avec risque factoriel, un plafond de
0.71~\%)~: les réseaux sont entraînés sur les mois 0 à 199, avec \omterm{def:ml:trees-and-boosting:early}{arrêt précoce} sur les mois 201 à 239 (règle du chapitre~5, un
mois purgé), et évalués sur les dix années de test avec le rapport de \texttt{firm.mlbase}, exactement comme les modèles du chapitre~4.
L'\omterm{def:ml:trees-and-boosting:early}{arrêt précoce} les interrompt après deux à quatre \omterm{def:ml:neural-networks-for-noisy-tabular-data:backprop}{époques} (\cref{fig:ml:nets:history})~: sur des données aussi bruitées, le
score de validation culmine presque immédiatement, et chaque \omterm{def:ml:neural-networks-for-noisy-tabular-data:backprop}{époque} supplémentaire ajuste davantage de bruit que de signal.

\begin{omfigure}
\begin{tikzpicture}
\begin{axis}[width=10.5cm,height=5.2cm,xlabel={époque},ylabel={$R^2$ de validation (\%)},tick label style={font=\small},
  label style={font=\small},xmin=0.5,xmax=9.5,ymin=0.3,ymax=0.85,xtick={1,2,3,4,5,6,7,8,9},grid=major,
  grid style={gray!25},legend columns=3,legend style={font=\footnotesize,at={(0.5,-0.3)},anchor=north,draw=none,
  fill=none,/tikz/every even column/.append style={column sep=8pt}},table/col sep=comma,unbounded coords=jump]
\addplot[omDef,very thick,mark=*] table[x=epoch,y=s1]{figdata/ml/07-neural-networks-for-noisy-tabular-data/history.csv};
\addplot[omThm,very thick,mark=square*] table[x=epoch,y=s2]{figdata/ml/07-neural-networks-for-noisy-tabular-data/history.csv};
\addplot[omProp,very thick,mark=triangle*] table[x=epoch,y=s3]{figdata/ml/07-neural-networks-for-noisy-tabular-data/history.csv};
\legend{graine 1,graine 2,graine 3}
\end{axis}
\end{tikzpicture}

\omcaption{R-carré de validation du petit réseau après chaque \omterm{def:ml:neural-networks-for-noisy-tabular-data:backprop}{époque}, pour trois graines~; l'entraînement s'arrête cinq \omterm{def:ml:neural-networks-for-noisy-tabular-data:backprop}{époques} après le
meilleur point. Les graines culminent aux \omterm{def:ml:neural-networks-for-noisy-tabular-data:backprop}{époques} 2, 2 et 4. Données~: \texttt{ml\_nets.fitted}.}\label{fig:ml:nets:history}
\end{omfigure}

Avec trois graines par variante, un \omterm{def:ml:neural-networks-for-noisy-tabular-data:dropout}{dropout} de 0.2 fait passer le IC moyen de test de 0.059 à 0.062, une \omterm{def:ml:neural-networks-for-noisy-tabular-data:adam}{décroissance des poids} de 0.01 ne change rien
à cette échelle, et les deux normalisations nuisent~: \omterm{def:ml:neural-networks-for-noisy-tabular-data:dropout}{normalisation par lot} à 0.052, \omterm{def:ml:neural-networks-for-noisy-tabular-data:dropout}{normalisation par couche} à 0.056. Normaliser
les unités cachées aident les réseaux profonds à s'entraîner~; ceux-ci sont peu profonds, leurs entrées sont déjà des rangs, et les statistiques d'un lot sont
une source de bruit de plus. Les écarts se chiffrent en quelques millièmes de IC, comparables au bruit des graines mesuré ensuite, et
trois graines par variante sont trop peu nombreuses pour les classer avec confiance.

\section{Les graines comme source de variance}

Dix graines du petit réseau (32 et 16 unités) donnent des IC de test allant de 0.057 à 0.065, avec une moyenne de 0.0601 et un
écart type de 0.0023~; dix graines du réseau plus grand (64, 32 et 16) donnent en moyenne 0.0619 (\cref{fig:ml:nets:seeds}).
Leur R-carré varie bien plus que leur IC, de 0.08~\% à 0.30~\% pour le petit réseau~: la graine modifie l'échelle des
prévisions, ce que le R-carré pénalise et ce qu'un classement ignore (chapitre~1). Le ratio de Sharpe net du portefeuille par déciles passe
de 1.37 à 1.78.

\begin{proposition}[Combien de graines une comparaison exige]\label{prop:ml:nets:seeds}
Si les scores de deux architectures varient selon les graines avec un écart type $s$, de façon indépendante, et si leurs moyennes diffèrent de
$\delta$, la différence de leurs moyennes sur $n$ graines a une erreur type $s\sqrt{2/n}$, donc elle atteint deux erreurs types
lorsque $n \ge 8s^2/\delta^2$.
\end{proposition}

\begin{proof}
$\Var(\bar a - \bar b) = s^2/n + s^2/n$~; résoudre $\delta = 2s\sqrt{2/n}$ pour $n$.
\end{proof}

Avec $s = 0.0024$ (écart type groupé) et $\delta = 0.0019$, la comparaison requiert 14 graines par architecture. Une seule exécution de chacune,
la pratique habituelle, qui compare deux tirages dont la différence a une erreur type de 0.0035, soit presque le double de l'effet.

\begin{definition}[Ensemble profond]\label{def:ml:neural-networks-for-noisy-tabular-data:ensemble}
Un \emph{ensemble profond}\index{ensemble profond} (deep ensemble) moyenne les prédictions de réseaux de même architecture entraînés à partir de
graines aléatoires différentes sur les mêmes données (Lakshminarayanan, Pritzel et Blundell, 2017).
\end{definition}

Faire la moyenne élimine la variance que les graines introduisent. L'ensemble à dix graines du petit réseau atteint un IC de 0.066, au-dessus
de chacun de ses membres, avec un R-carré de 0.32~\% et un ratio de Sharpe net de 1.73~; l'ensemble du réseau plus grand
atteint 0.066 également. La question de l'architecture qui opposait les deux chercheurs disparaît dès que chaque réseau
est moyenné sur ses graines~: la moyenne a apporté ce que l'architecture ne pouvait pas. C'est à nouveau la \cref{prop:ml:trees:bagging}
avec les graines jouant le rôle des échantillons bootstrap.

\begin{omfigure}
\begin{tikzpicture}
\begin{axis}[width=9cm,height=5.6cm,ylabel={coefficient d'information de test},xmin=0.4,xmax=2.6,ymin=0.054,ymax=0.068,
  xtick={1,2},xticklabels={{32, 16 unités},{64, 32, 16 unités}},tick label style={font=\small},label style={font=\small},
  grid=major,grid style={gray!25},yticklabel style={/pgf/number format/fixed,/pgf/number format/precision=3},
  scaled y ticks=false,legend columns=3,legend style={font=\footnotesize,at={(0.5,-0.2)},anchor=north,draw=none,
  fill=none,/tikz/every even column/.append style={column sep=8pt}},table/col sep=comma]
\addplot[only marks,mark=*,mark size=1.8pt,gray] table[x=arch,y=ic]{figdata/ml/07-neural-networks-for-noisy-tabular-data/seeds.csv};
\addplot[only marks,mark=-,mark size=9pt,very thick,omDef] table[x=arch,y=mean]{figdata/ml/07-neural-networks-for-noisy-tabular-data/means.csv};
\addplot[only marks,mark=star,mark size=4pt,thick,omThm] table[x=arch,y=ensemble]{figdata/ml/07-neural-networks-for-noisy-tabular-data/means.csv};
\legend{une graine,moyenne de 10 graines,ensemble de 10}
\end{axis}
\end{tikzpicture}

\omcaption{IC de test de dix graines de deux architectures sur le panel du chapitre~4, leurs moyennes, et l'IC de chaque
ensemble de dix graines. Données~: \texttt{ml\_nets.seed\_study} et \texttt{ml\_nets.ensemble}.}\label{fig:ml:nets:seeds}
\end{omfigure}

\section{Ensembles}

Un ensemble de graines coûte $n$ entraînements et rien d'autre, et il se parallélise sans effort. Il apporte aussi quelque chose qu'un
seul réseau ne peut fournir~: la dispersion des prévisions de ses membres pour chaque ligne, que le chapitre~11 utilise comme mesure de
l'incertitude propre au modèle. Les ensembles d'architectures différentes, ou de réseaux associés à des arbres boostés et à une
régression ridge (stacking, livre~7, chapitre~14), vont plus loin lorsque les erreurs des membres sont moins corrélées que celles des graines.

\begin{method}[Entraîner un réseau sur des données tabulaires bruitées]\label{met:ml:nets:training}
\begin{enumerate}
\item Trier les variables par date~; diviser la cible par son écart type d'entraînement.
\item Commencer petit (deux couches cachées, quelques dizaines d'unités), avec AdamW, un \omterm{def:ml:trees-and-boosting:boosting}{taux d'apprentissage} voisin de $10^{-3}$, et un \omterm{def:ml:trees-and-boosting:early}{arrêt précoce} sur un
bloc de validation purgé.
\item Figer chaque source d'aléa (initialisation, \omterm{def:ml:neural-networks-for-noisy-tabular-data:dropout}{dropout}, ordre des lots) à une graine par exécution~; entraîner plusieurs graines de
tout ce qu'on compare, et présenter la moyenne et la dispersion.
\item Déployer l'ensemble de graines, et non la meilleure graine.
\item Comparer avec la régression ridge et les arbres boostés, sur les mêmes mois et avec le même rapport.
\end{enumerate}
\end{method}

\section{Réseaux contre boosting sur des données tabulaires}

La comparaison avec le chapitre~4 se fait à conditions égales~: même panel, mêmes années de test, même rapport. Un seul petit réseau
(IC 0.060, ratio de Sharpe net de 1.37--1.78 selon les graines) égale la régression ridge (0.060, 1.59)~; son ensemble de dix graines
(0.066, 1.73) égale l'IC du lasso et reste derrière les arbres boostés par défaut de la maison (0.075, 2.02). Gu, Kelly et Xiu ont trouvé
des réseaux légèrement devant les arbres sur des actions américaines réelles, avec soixante ans de données, et Grinsztajn, Oyallon et Varoquaux des arbres
devant les réseaux sur la plupart des bancs d'essai tabulaires~: sur des panels de variables classées, les deux familles sont proches, le
réseau a besoin d'un ensemble pour être compétitif, et l'ensemble d'arbres coûte moins cher à régler.

\section{Tutoriel~: vingt graines}\label{tut:ml:neural-networks-for-noisy-tabular-data:tut}

\begin{tutorial}
\textbf{Objectif.} Entraîner de petits réseaux et des réseaux plus grands à dix graines chacun, les regrouper en ensemble, essayer quatre régularisateurs et comparer
à la régression ridge et aux arbres boostés. \textbf{État final~:} les \cref{fig:ml:nets:history,fig:ml:nets:seeds}, les chiffres du chapitre.
\begin{tutsteps}
\item \textbf{La boucle d'entraînement déterministe}~: un générateur pour l'ordre des lots, la graine pour l'initialisation et
le \omterm{def:ml:neural-networks-for-noisy-tabular-data:dropout}{dropout}, un \omterm{def:ml:trees-and-boosting:early}{arrêt précoce} fondé sur la perte de validation.
\omcode{code/firm/nettab/firm_nettab.py}{77}{115}{Boucle d'entraînement déterministe avec arrêt précoce.}\label{lst:ml:nets:fit}
\item \textbf{Exécuter} \texttt{ml\_nets.seed\_study()}, \texttt{ensemble()}, \texttt{ablation()}, \texttt{baselines()} et
\texttt{fig\_nets.py}.
\end{tutsteps}
\textbf{Pour aller plus loin.} Entraîner avec 30 graines par architecture et vérifier la prédiction de la \cref{prop:ml:nets:seeds}~;
remplacer ReLU par tanh et voir si la dispersion des graines change.
\end{tutorial}

\section{Construction~: les réseaux tabulaires de la maison}\label{bld:ml:neural-networks-for-noisy-tabular-data:nettab}

\begin{build}
\bfield{Rôle} Tout réseau que la maison entraîne sur des données tabulaires est reproductible bit à bit et comparé d'une graine à l'autre.
\bfield{Interface} \texttt{set\_determinism(threads)}, \texttt{MLP(p, hidden, dropout, norm, activation)},
\texttt{fit(X, y, Xv, yv, hidden, dropout, norm, lr, weight\_decay, batch, epochs, patience, seed) -> Fitted} avec
\texttt{predict}, \texttt{history}, \texttt{best\_epoch}, \texttt{state\_hash()}~; \texttt{SeedEnsemble(seeds, **kw)}.
\bfield{Règles} Un seul thread, algorithmes déterministes, chaque tirage aléatoire issu de la graine de l'exécution~; entrées standardisées avec
les statistiques d'entraînement stockées dans le modèle~; le bloc de validation purgé par l'appelant.
\bfield{Tests d'acceptation} \texttt{code/firm/nettab/tests/}~: deux exécutions avec une même graine donnent des empreintes de poids identiques, deux
graines donnent des empreintes différentes~; le réseau apprend une fonction non linéaire injectée~; l'\omterm{def:ml:trees-and-boosting:early}{arrêt précoce} conserve les poids de la meilleure
\omterm{def:ml:neural-networks-for-noisy-tabular-data:backprop}{époque}~; l'ensemble fait la moyenne de ses membres.
\bfield{Pour aller plus loin} Entraînement en précision mixte et chargeur de données pour des panels plus grands (chapitre~23)~; exportation des poids pour
l'inférence quantifiée du chapitre~26.
\end{build}

\begin{omsources}
\item D.~P.~Kingma et J.~Ba, ``Adam: a method for stochastic optimization'', \emph{ICLR}, 2015.
\item I.~Loshchilov et F.~Hutter, ``Decoupled weight decay regularization'', \emph{ICLR}, 2019.
\item N.~Srivastava, G.~Hinton, A.~Krizhevsky, I.~Sutskever et R.~Salakhutdinov, ``Dropout: a simple way to prevent
neural networks from overfitting'', \emph{JMLR} 15, 2014.
\item S.~Ioffe et C.~Szegedy, ``Batch normalization'', \emph{ICML}, 2015~; J.~L.~Ba, J.~R.~Kiros et G.~E.~Hinton,
``Layer normalization'', arXiv, 2016.
\item B.~Lakshminarayanan, A.~Pritzel et C.~Blundell, ``Simple and scalable predictive uncertainty estimation using deep
ensembles'', \emph{NeurIPS}, 2017.
\item D.~E.~Rumelhart, G.~E.~Hinton et R.~J.~Williams, ``Learning representations by back-propagating errors'',
\emph{Nature} 323, 1986.
\end{omsources}

\section{Exercices}

\begin{exercise}[$\star$]\label{exo:ml:neural-networks-for-noisy-tabular-data:1}
Compter les paramètres d'un \omterm{def:ml:neural-networks-for-noisy-tabular-data:mlp}{perceptron multicouche} à 40 entrées, des couches cachées de 32 et 16 unités et une sortie. Et avec une première couche
de 64 unités ajoutée~?
\end{exercise}

\begin{exercise}[$\star$]\label{exo:ml:neural-networks-for-noisy-tabular-data:2}
Les ICs de test de deux architectures varient selon les graines avec un écart type de 0.003 et diffèrent de 0.001 en moyenne. Combien
de graines par architecture une comparaison à deux erreurs types exige-t-elle~?
\end{exercise}

\begin{exercise}[$\star$]\label{exo:ml:neural-networks-for-noisy-tabular-data:3}
Avec 160\,000 lignes d'entraînement et des \omterm{def:ml:neural-networks-for-noisy-tabular-data:backprop}{mini-lots} de 512, combien de pas de gradient compte une \omterm{def:ml:neural-networks-for-noisy-tabular-data:backprop}{époque}~? Et pour les
100\,000 lignes du chapitre~?
\end{exercise}

\begin{exercise}[$\star\star$]\label{exo:ml:neural-networks-for-noisy-tabular-data:4}
Pourquoi la graine fait-elle bouger le R-deux (de 0.08~\% à 0.30~\%) bien davantage que l'IC (de 0.057 à 0.065)~?
\end{exercise}

\begin{exercise}[$\star\star$]\label{exo:ml:neural-networks-for-noisy-tabular-data:5}
Pourquoi la \omterm{def:ml:neural-networks-for-noisy-tabular-data:dropout}{normalisation par lot} peut-elle nuire à un réseau peu profond sur des entrées rangées, et que faut-il attendre d'un réseau profond~?
\end{exercise}

\begin{exercise}[$\star\star$]\label{exo:ml:neural-networks-for-noisy-tabular-data:6}
\emph{Trouver la faille.} «~Nous avons entraîné vingt graines et déployé celle qui obtenait le meilleur score de validation.~»
\end{exercise}

\begin{exercise}[$\star\star\star$]\label{exo:ml:neural-networks-for-noisy-tabular-data:7}
\emph{Programmation.} Calculer l'IC des ensembles formés des 2, 5 et 10 premières graines du petit réseau. Quelle part du gain entre
une et dix graines les deux premières procurent-elles déjà~?
\end{exercise}

\begin{exercise}[$\star\star\star$]\label{exo:ml:neural-networks-for-noisy-tabular-data:8}
Montrer que, pour une couche ReLU et une perte quadratique, le gradient de la \omterm{def:ml:neural-networks-for-noisy-tabular-data:backprop}{rétropropagation} par rapport à $A_1$ vaut $\sum_i
(\partial\ell_i/\partial z_{1,i})\odot\mathbf 1\{A_1x_i + c_1 > 0\}\,x_i^\top$, puis comparer son coût à celui de la
passe avant.
\end{exercise}

\section{Problème~: Vingt graines}

\begin{problem}[{Problème du week-end --- un réseau comme variable aléatoire}]\label{pb:ml:neural-networks-for-noisy-tabular-data:1}
Le panel du chapitre~4~; réseaux entraînés sur les mois 0 à 199, arrêtés sur les mois 201 à 239, évalués sur les mois 240 à 359.

\medskip\noindent\textbf{Partie I --- Entraînement.}
\begin{enumerate}
\item Combien de paramètres possèdent les deux réseaux~?
\item Après combien d'\omterm{def:ml:neural-networks-for-noisy-tabular-data:backprop}{époques} l'\omterm{def:ml:trees-and-boosting:early}{arrêt précoce} (early stopping) les arrête-t-il, et pourquoi si tôt~?
\item Que font le \omterm{def:ml:neural-networks-for-noisy-tabular-data:dropout}{dropout}, la \omterm{def:ml:neural-networks-for-noisy-tabular-data:adam}{décroissance des poids} et les deux normalisations à l'IC~?
\item Pourquoi trois graines par variante sont-elles insuffisantes pour les classer~?
\end{enumerate}

\medskip\noindent\textbf{Partie II --- Graines.}
\begin{enumerate}[resume]
\item Quelle plage d'IC, de R-deux et de ratio de Sharpe net couvrent dix graines du petit réseau~?
\item Quels sont les ICs moyens des deux architectures, leur écart et l'écart type poolé des graines~?
\item Combien de graines par architecture une comparaison à deux erreurs types exige-t-elle~?
\item Quelle est l'erreur type d'une comparaison d'une exécution de chaque modèle~?
\end{enumerate}

\medskip\noindent\textbf{Partie III --- Ensembles.}
\begin{enumerate}[resume]
\item Quels scores obtiennent les deux ensembles de dix graines~?
\item Pourquoi l'ensemble bat-il chacun de ses membres~?
\item Que devient la question de l'architecture une fois chaque réseau mis en ensemble~?
\item Qu'apporte un ensemble au-delà d'une meilleure prévision~?
\end{enumerate}

\medskip\noindent\textbf{Partie IV --- Le verdict.}
\begin{enumerate}[resume]
\item Comment les réseaux se comparent-ils à la régression ridge, au lasso et aux arbres boostés sur les mêmes mois~?
\item Énoncer le \emph{résultat nommé}~: l'écart type des graines de l'IC de test face à l'écart d'architecture, et le
nombre de graines nécessaire pour le résoudre à deux erreurs types.
\item Quel modèle faudrait-il déployer, et sous quelle forme~?
\item Que doit consigner un journal de recherche à propos d'un réseau~?
\item Comment le tableau changerait-il avec soixante ans de données~?
\item Que garantit le hachage d'état, et que ne garantit-il pas~?
\item D'où provient l'aléa de l'entraînement d'un réseau~?
\item En une phrase~: qu'est-ce qu'un réseau entraîné une seule fois sur des données bruitées~?
\end{enumerate}
\end{problem}

\section{Questions d'entretien}

\begin{interviewq}[$\star$ \iqroles{ingénieur ML}]\label{iq:ml:neural-networks-for-noisy-tabular-data:1}
Que calcule la \omterm{def:ml:neural-networks-for-noisy-tabular-data:backprop}{rétropropagation}, et pourquoi son coût est-il faible~?
\end{interviewq}

\begin{interviewq}[$\star$ \iqroles{ingénieur ML, chercheur}]\label{iq:ml:neural-networks-for-noisy-tabular-data:2}
Expliquez la différence entre \omterm{def:ml:neural-networks-for-noisy-tabular-data:adam}{Adam} et AdamW.
\end{interviewq}

\begin{interviewq}[$\star\star$ \iqroles{chercheur, ingénieur ML}]\label{iq:ml:neural-networks-for-noisy-tabular-data:3}
Votre nouvelle architecture bat l'ancienne de 0.003 d'IC en une exécution chacune. Que faites-vous avant d'y croire~?
\end{interviewq}

\begin{interviewq}[$\star\star$ \iqroles{ingénieur ML}]\label{iq:ml:neural-networks-for-noisy-tabular-data:4}
Comment rendez-vous un entraînement PyTorch reproductible sur CPU~? Qu'est-ce qui peut encore casser la reproductibilité~?
\end{interviewq}

\begin{interviewq}[$\star\star$ \iqroles{chercheur}]\label{iq:ml:neural-networks-for-noisy-tabular-data:5}
\omterm{def:ml:neural-networks-for-noisy-tabular-data:mlp}{Réseau de neurones} ou \omterm{def:ml:trees-and-boosting:boosting}{gradient boosting} pour un modèle de rendements en coupe transversale~? Argumentez les deux positions.
\end{interviewq}

\begin{interviewq}[$\star\star\star$ \iqroles{chercheur, ingénieur ML}]\label{iq:ml:neural-networks-for-noisy-tabular-data:6}
Pourquoi moyenner des réseaux entraînés à partir de graines différentes améliore-t-il la prévision, et quand cesse-t-il d'aider~?
\end{interviewq}
